
%%%%%%%%%%%%%%%%%%%%%%%%%%%%%%%%%%%%%%%%%%
%%          Resultados Obtidos          %%
%%%%%%%%%%%%%%%%%%%%%%%%%%%%%%%%%%%%%%%%%%

% Neste capítulo são expostos os resultados de sua pesquisa.
% No caso do TCC I, resultados esperados ou parciais, para TCC II, resultados "finais".

\chapter{Resultados}
    \label{cha:resultados}
    % remove nome do capítulo do cabeçalho

    Com a proposta elaborada neste trabalho, é esperado a obtenção de uma plataforma capaz de integrar facilmente os dados; que seja possível estender as funcionalidades para vários algoritmos de aprendizado de máquina; além claro, da extensão da própria interface, tudo isso para que a plataforma atenda a demanda de diversos tipos de problemas.

    Em relação aos dados, a plataforma deverá oferecer a suporte a diferentes fontes de arquivos, tanto arquivos locais, quanto arquivos em nuvem, como, por exemplo, o Google Drive. Além disso, a plataforma deverá conseguir integrar dados de diferentes tipos, como, por exemplo, dados de um banco de dados relacional, dados de um aquivo de planilha como XLXS e dados de um arquivo CSV.

    A plataforma também deverá integrar novos algoritmos a partir das interfaces herdadas do \textit{scikit-learn}, como, por exemplo, o \textit{Random Forest} e o \textit{Gradient Boosting}. Além disso, a plataforma deverá ser capaz de integrar novos algoritmos de aprendizado de máquina, como, por exemplo, o \textit{Support Vector Machine} e o \textit{Neural Network}, com a condição da utilização das interfaces compatíveis. Desta forma será possível integrar novos algoritmos com poucas linhas de código.

    A interface, no que lhe concerne, deverá ser capaz integrar novas funcionalidades, como, por exemplo, a criação de novos tipos de gráficos, como, por exemplo, gráficos de dispersão, gráficos de barras e gráficos de pizza. Além disso, a interface deverá conseguir integrar novas funcionalidades, como, por exemplo, a integração de novos algoritmos com parâmetros personalizados, como, por exemplo, o algoritmo \textit{KNN} com o parâmetro \textit{weights}.

    Desta forma, a plataforma deverá ser capaz de atender a demanda de diversos tipos de problemas, como, por exemplo, problemas de classificação, problemas de regressão e problemas de agrupamento. Seguindo os requisitos propostos, a plataforma deverá conseguir atender a demanda levantada no começo do trabalho com os pesquisadores do \textit{SmartEnergy}.
